% Written by scripts/87_supplement_tables.py. Do not edit by hand.
% Receipts: analysis/brm/95_null_gap_equivalence.csv
\begingroup
\singlespacing
\footnotesize
\setlength{\tabcolsep}{4pt}
\begin{longtable}{>{\raggedright\arraybackslash}p{3.4cm}lll}
\caption{Level 2, Exploratory: Equivalence of Simulated Gaps Where the Population Shows None}\label{tab:s4null}\\
\toprule
Dataset & $\delta = 0.10$ SD & $\delta = 0.20$ SD & $\delta = 0.25$ SD \\
\midrule
\endfirsthead
\multicolumn{4}{@{}l}{\textit{Table \thetable, continued}} \\
\toprule
Dataset & $\delta = 0.10$ SD & $\delta = 0.20$ SD & $\delta = 0.25$ SD \\
\midrule
\endhead
\bottomrule
\endlastfoot
Worked example (4) & 3/0/1 & 4/0/0 & 4/0/0 \\
Bisbee et al. (23) & 5/11/7 & 13/4/6 & 18/3/2 \\
Argyle et al. (23) & 7/3/13 & 13/1/9 & 14/1/8 \\
\end{longtable}
\vspace{-6pt}
\noindent\textit{Note.} Counts are pass / fail / unresolved for the contrasts stopped as no population gap. A contrast passes when the 90\% interval of the simulated gap lies inside $(-\delta, \delta)$ and fails when it lies wholly outside; $\delta$ is in reference standard deviations. Not part of the frozen ladder. Every row is in \texttt{analysis/brm/95\_null\_gap\_equivalence.csv}.
\par
\endgroup
